\documentclass[12pt,a4paper]{article}

\usepackage[utf8]{inputenc}
\usepackage[T1]{fontenc}
\usepackage[spanish]{babel}
\usepackage{geometry}
\usepackage{graphicx}
\usepackage{xcolor}
\usepackage{booktabs}
\usepackage{tabularx}
\usepackage{longtable}
\usepackage{enumitem}
\usepackage{listings}
\usepackage{hyperref}
\usepackage{fancyhdr}
\usepackage{titlesec}
\usepackage{tcolorbox}
\usepackage{multirow}
\usepackage{array}
\usepackage{float}

\geometry{margin=2.5cm}

\definecolor{codegreen}{rgb}{0.0, 0.5, 0.0}
\definecolor{codegray}{rgb}{0.5, 0.5, 0.5}
\definecolor{codepurple}{rgb}{0.58, 0.0, 0.82}
\definecolor{backcolour}{rgb}{0.95, 0.95, 0.95}
\definecolor{gradeA}{HTML}{2E7D32}
\definecolor{gradeB}{HTML}{1565C0}
\definecolor{gradeC}{HTML}{E65100}
\definecolor{accentgreen}{HTML}{4CAF50}
\definecolor{accentred}{HTML}{F44336}
\definecolor{accentorange}{HTML}{FF9800}
\definecolor{sectionblue}{HTML}{1A237E}

\lstdefinestyle{jsstyle}{
    backgroundcolor=\color{backcolour},
    commentstyle=\color{codegreen},
    keywordstyle=\color{codepurple}\bfseries,
    numberstyle=\tiny\color{codegray},
    stringstyle=\color{codegreen},
    basicstyle=\ttfamily\footnotesize,
    breakatwhitespace=false,
    breaklines=true,
    captionpos=b,
    keepspaces=true,
    numbers=left,
    numbersep=5pt,
    showspaces=false,
    showstringspaces=false,
    showtabs=false,
    tabsize=2,
    frame=single,
    rulecolor=\color{codegray},
    language=Java,
    morekeywords={const, let, async, await, class, constructor, try, catch, finally, return, if, else, for, of, new, this, function, require, module, exports, process, setTimeout, setInterval, clearTimeout, clearInterval}
}
\lstset{style=jsstyle}

\pagestyle{fancy}
\fancyhf{}
\fancyhead[L]{\small\textsc{Reporte de Implementaci\'on --- Stack Windows Electron}}
\fancyhead[R]{\small\thepage}
\fancyfoot[C]{\small\textit{Documento generado el 17 de marzo de 2026}}
\renewcommand{\headrulewidth}{0.4pt}
\renewcommand{\footrulewidth}{0.2pt}

\titleformat{\section}{\Large\bfseries\color{sectionblue}}{\thesection}{1em}{}
\titleformat{\subsection}{\large\bfseries\color{sectionblue!80}}{\thesubsection}{1em}{}
\titleformat{\subsubsection}{\normalsize\bfseries\color{sectionblue!60}}{\thesubsubsection}{1em}{}

\newtcolorbox{findingbox}[1][]{
    colback=red!5!white,
    colframe=accentred!75!black,
    fonttitle=\bfseries,
    title=#1
}

\newtcolorbox{fixbox}[1][]{
    colback=green!5!white,
    colframe=accentgreen!75!black,
    fonttitle=\bfseries,
    title=#1
}

\newtcolorbox{infobox}[1][]{
    colback=blue!5!white,
    colframe=sectionblue!75!black,
    fonttitle=\bfseries,
    title=#1
}

\newtcolorbox{regressionbox}[1][]{
    colback=orange!5!white,
    colframe=accentorange!75!black,
    fonttitle=\bfseries,
    title=#1
}

\begin{document}

% === PORTADA ===
\begin{titlepage}
\centering
\vspace*{3cm}
{\Huge\bfseries\color{sectionblue} Reporte de Implementaci\'on\\[0.5cm]}
{\LARGE\color{sectionblue!70} Plan Refinado de Mejoras\\[0.3cm]}
{\large\color{codegray} Stack Windows Electron --- Organizador Espacial de Ventanas}
\vspace{2cm}
\begin{tcolorbox}[colback=sectionblue!5, colframe=sectionblue, width=12cm]
\centering\large
\textbf{%%TOTAL_TASKS%% tareas completadas} $\cdot$ \textbf{3 fases ejecutadas}\\[0.3cm]
\textbf{%%TOTAL_TESTS%% tests automatizados} $\cdot$ \textbf{%%TOTAL_COMMITS%% commits}\\[0.3cm]
\textbf{Calificaci\'on: B+ $\rightarrow$ A--}
\end{tcolorbox}
\vspace{2cm}
\begin{tabular}{ll}
\textbf{Proyecto:} & Stack Windows Electron \\
\textbf{Versi\'on:} & 1.0.0 \\
\textbf{Rama:} & \texttt{fix/dimensiones-limitadas} \\
\textbf{Fecha:} & 17 de marzo de 2026 \\
\textbf{Plan base:} & Refined Plan (9 revisiones de expertos) \\
\textbf{Plataforma:} & Windows (exclusivo) \\
\end{tabular}
\vfill
{\small Reporte generado como parte del proceso de mejora continua del proyecto.}
\end{titlepage}

\tableofcontents
\newpage

% === SECTION 1: RESUMEN EJECUTIVO ===
\section{Resumen Ejecutivo}

Se ejecut\'o el plan refinado de mejoras para \textbf{Stack Windows Electron}, un plan de %%TOTAL_TASKS%% tareas distribuidas en 3 fases que abord\'o los hallazgos pendientes de la auditor\'ia integral previa (calificaci\'on B+). El plan fue dise\~nado a partir de 9 revisiones de expertos que identificaron riesgos cr\'iticos en el ordenamiento de tareas, modelos de concurrencia defectuosos y una estrategia de testing incompleta.

Como resultado, se implementaron mejoras fundamentales en estabilidad (file locking con \texttt{proper-lockfile}, arquitectura CQS, heartbeat funcional), seguridad (eliminaci\'on de \texttt{unsafe-inline} del CSP, endurecimiento de \texttt{webPreferences}), calidad (ESLint + Prettier, separaci\'on del renderer monol\'itico), testing (%%TOTAL_TESTS%% tests automatizados cubriendo unit, integraci\'on y UI), y rendimiento (cache de callbacks FFI, DPI awareness, deduplicaci\'on de l\'ogica). La calificaci\'on estimada del proyecto mejor\'o de \textbf{B+} a \textbf{A--}.

% === SECTION 2: ORIGEN Y CONTEXTO ===
\section{Origen y Contexto}

\subsection{Auditor\'ia Previa}
La auditor\'ia integral realizada el 5 de marzo de 2026 evalu\'o 10 categor\'ias sobre 9 archivos fuente (3,168 l\'ineas). De los 17 hallazgos identificados, 11 fueron corregidos durante la auditor\'ia y 6 quedaron pendientes como recomendaciones futuras.

\subsection{Hallazgos Pendientes y su Resoluci\'on}

\begin{table}[H]
\centering
\renewcommand{\arraystretch}{1.3}
\begin{tabularx}{\textwidth}{c X c}
\toprule
\textbf{\#} & \textbf{Recomendaci\'on Pendiente} & \textbf{Estado Actual} \\
\midrule
5 & Separar \texttt{index.html} en archivos HTML/CSS/JS independientes & \textcolor{accentgreen}{\textbf{COMPLETADO}} \\
7 & Agregar file locking al registro de instancias & \textcolor{accentgreen}{\textbf{COMPLETADO}} \\
9 & Cachear callback koffi en \texttt{getAvailableWindows} & \textcolor{accentgreen}{\textbf{COMPLETADO}} \\
10 & Agregar infraestructura de testing b\'asica & \textcolor{accentgreen}{\textbf{COMPLETADO}} \\
12 & Agregar eslint + prettier & \textcolor{accentgreen}{\textbf{COMPLETADO}} \\
\bottomrule
\end{tabularx}
\caption{Estado de las recomendaciones pendientes de la auditor\'ia}
\end{table}

\subsection{El Plan Refinado}
El plan original fue sometido a 9 revisiones de expertos que identificaron: (1) riesgo de corrupci\'on de datos si el file locking no se implementaba primero, (2) un modelo de caching incompatible con el locking propuesto, (3) estimaciones de esfuerzo subestimadas, y (4) ausencia de tests de integraci\'on y E2E. El plan refinado adopt\'o un enfoque ``safety-first'' con arquitectura CQS y un patr\'on Lock-Modify-Write estricto.

% === SECTION 3: METODOLOGIA ===
\section{Metodolog\'ia de Ejecuci\'on}

\subsection{Modelo de Trabajo}
Se utiliz\'o un modelo de \textbf{swarm} con workers aut\'onomos ejecutando tareas en paralelo. Cada worker oper\'o en un \texttt{git worktree} aislado, garantizando que los cambios no interfirieran entre s\'i durante el desarrollo.

\subsection{Flujo por Fases}
\begin{enumerate}
    \item \textbf{Crear worktrees} para las tareas sin dependencias
    \item \textbf{Despachar workers} en paralelo (3--6 simult\'aneos)
    \item \textbf{Integrar} resultados en la rama principal con \texttt{git merge --no-ff}
    \item \textbf{Verificar} que \texttt{npm test} y \texttt{npm run lint} pasan
    \item \textbf{Repetir} con la siguiente oleada de tareas
\end{enumerate}

\subsection{Oleadas de Ejecuci\'on}

\begin{table}[H]
\centering
\renewcommand{\arraystretch}{1.3}
\begin{tabularx}{\textwidth}{l l X}
\toprule
\textbf{Oleada} & \textbf{Tareas} & \textbf{Descripci\'on} \\
\midrule
Fase 1, Ola 1 & z2b.1, z2b.3, z2b.6, z2b.8 & File locking, separar renderer, ESLint, WebPreferences \\
Fase 1, Ola 2 & z2b.2, z2b.5, z2b.7 & CQS+heartbeat, eliminar unsafe-inline, fix lint \\
Fase 2, Ola 1 & z2b.9 & Configurar Vitest \\
Fase 2, Ola 2 & z2b.11--z2b.15 & Tests unitarios e integraci\'on \\
Fase 2+3 & z2b.10, z2b.14, z2b.16--z2b.19 & Tests pendientes + robustez y rendimiento \\
\bottomrule
\end{tabularx}
\caption{Oleadas de ejecuci\'on del plan}
\end{table}

% === SECTION 4: FASE 1 ===
\section{Fase 1 --- Estabilidad y Calidad Fundacional}

\subsection{File Locking para InstanceRegistry (z2b.1)}

\begin{findingbox}[Problema: Race condition en acceso concurrente]
El archivo \texttt{instance-registry.json} era le\'ido y escrito por m\'ultiples instancias de la aplicaci\'on sin ning\'un mecanismo de locking. Dos instancias escribiendo simult\'aneamente pod\'ian corromper el archivo JSON.
\end{findingbox}

\begin{fixbox}[Soluci\'on: proper-lockfile con Lock-Modify-Write]
Se integr\'o la librer\'ia \texttt{proper-lockfile} con configuraci\'on robusta: \texttt{stale: 10000ms} (locks abandonados se liberan autom\'aticamente), \texttt{retries: 3} con backoff exponencial. Todas las operaciones de escritura siguen el patr\'on Lock $\rightarrow$ Modify $\rightarrow$ Write $\rightarrow$ Unlock en un bloque \texttt{try/finally}.
\end{fixbox}

\subsection{Arquitectura CQS y Heartbeat (z2b.2)}

\begin{findingbox}[Problema: Side effects ocultos y heartbeat inexistente]
\texttt{getOtherInstancesHwnds()} le\'ia, podaba Y escrib\'ia el registro --- un efecto secundario oculto en una operaci\'on de lectura. Adem\'as, \texttt{\_startHeartbeat()} se invocaba en \texttt{init()} pero \textbf{el m\'etodo no exist\'ia}, causando un error silencioso.
\end{findingbox}

\begin{fixbox}[Soluci\'on: Command-Query Separation]
Se refactoriz\'o el m\'odulo para separar estrictamente queries (solo lectura) de commands (escritura):
\end{fixbox}

\begin{table}[H]
\centering
\renewcommand{\arraystretch}{1.3}
\begin{tabularx}{\textwidth}{l l X}
\toprule
\textbf{Tipo} & \textbf{M\'etodo} & \textbf{Comportamiento} \\
\midrule
Query & \texttt{getOtherInstancesHwnds()} & Solo lectura --- nunca escribe al disco \\
Query & \texttt{getRegistry()} & Solo lectura --- retorna datos actuales \\
Command & \texttt{\_register()} & Lock $\rightarrow$ Read $\rightarrow$ Modify $\rightarrow$ Write $\rightarrow$ Unlock \\
Command & \texttt{unregister()} & Lock $\rightarrow$ Read $\rightarrow$ Delete $\rightarrow$ Write $\rightarrow$ Unlock \\
Command & \texttt{\_startHeartbeat()} & Cada 60s: actualiza timestamp + poda instancias muertas \\
\bottomrule
\end{tabularx}
\caption{Separaci\'on de queries y commands en InstanceRegistry}
\end{table}

\subsection{Separaci\'on del Renderer Monol\'itico (z2b.3)}

\begin{findingbox}[Problema: Archivo monol\'itico de 1,715 l\'ineas]
\texttt{src/renderer/index.html} conten\'ia todo el HTML, CSS y JavaScript de la interfaz en un solo archivo, impidiendo linting, testing unitario y eliminaci\'on de \texttt{unsafe-inline} del CSP.
\end{findingbox}

\begin{fixbox}[Soluci\'on: Extracci\'on pura en 3 archivos]
Se realiz\'o una extracci\'on sin cambios de l\'ogica:
\end{fixbox}

\begin{table}[H]
\centering
\renewcommand{\arraystretch}{1.3}
\begin{tabular}{l r l r}
\toprule
\textbf{Antes} & \textbf{L\'ineas} & \textbf{Despu\'es} & \textbf{L\'ineas} \\
\midrule
\texttt{index.html} (todo) & 1,715 & \texttt{index.html} (solo HTML) & %%RENDERER_HTML_LINES%% \\
 & & \texttt{styles.css} & %%RENDERER_CSS_LINES%% \\
 & & \texttt{app.js} & %%RENDERER_JS_LINES%% \\
\bottomrule
\end{tabular}
\caption{Separaci\'on del renderer monol\'itico}
\end{table}

\subsection{Eliminaci\'on de unsafe-inline del CSP (z2b.5)}

\begin{findingbox}[Problema: CSP permit\'ia scripts inline]
El Content Security Policy inclu\'ia \texttt{script-src 'unsafe-inline'} porque todo el JavaScript estaba inline en el HTML. Los 11 atributos \texttt{onclick} tambi\'en constitu\'ian scripts inline.
\end{findingbox}

\begin{fixbox}[Soluci\'on: addEventListener + CSP endurecido]
Se reemplazaron los 11 atributos \texttt{onclick}/\texttt{onchange} con llamadas a \texttt{addEventListener} en \texttt{app.js}. El CSP actualizado:

\texttt{%%CSP_CURRENT%%}
\end{fixbox}

\subsection{ESLint y Prettier (z2b.6, z2b.7)}

Se configur\'o ESLint con reglas para Node.js/CommonJS y un override para archivos del renderer (entorno browser). Prettier se configur\'o con: single quotes, semicolons, trailing commas, 120 caracteres de ancho.

\begin{infobox}[Resultado del linting]
\textbf{%%LINT_ERRORS%% errores} $\cdot$ \textbf{%%LINT_WARNINGS%% warnings} (todos con justificaci\'on documentada via \texttt{eslint-disable} comments)
\end{infobox}

\subsection{Endurecimiento de WebPreferences (z2b.8)}

\begin{table}[H]
\centering
\renewcommand{\arraystretch}{1.3}
\begin{tabular}{l l l}
\toprule
\textbf{Propiedad} & \textbf{Valor} & \textbf{Estado} \\
\midrule
\texttt{contextIsolation} & \texttt{true} & Ya exist\'ia \\
\texttt{nodeIntegration} & \texttt{false} & Ya exist\'ia \\
\texttt{sandbox} & \texttt{true} & Ya exist\'ia \\
\texttt{webSecurity} & \texttt{true} & Ya exist\'ia \\
\texttt{webviewTag} & \texttt{false} & \textcolor{accentgreen}{\textbf{AGREGADO}} \\
\texttt{allowRunningInsecureContent} & \texttt{false} & \textcolor{accentgreen}{\textbf{AGREGADO}} \\
\bottomrule
\end{tabular}
\caption{Configuraci\'on de seguridad de WebPreferences}
\end{table}

% === SECTION 5: FASE 2 ===
\section{Fase 2 --- Infraestructura de Testing}

\subsection{Configuraci\'on de Vitest (z2b.9)}
Se configur\'o Vitest como test runner con mocks globales para \texttt{electron} y \texttt{koffi}. El principal desaf\'io t\'ecnico fue que Vitest v4 bloquea \texttt{require('vitest')} en m\'odulos CJS, requiriendo un archivo de setup en formato ESM (\texttt{.mjs}).

\subsection{Tests Unitarios por M\'odulo}

\begin{table}[H]
\centering
\renewcommand{\arraystretch}{1.3}
\begin{tabularx}{\textwidth}{l l r X}
\toprule
\textbf{M\'odulo} & \textbf{Archivo de Test} & \textbf{Tests} & \textbf{\'Areas Cubiertas} \\
\midrule
WindowManager & \texttt{window-manager.test.js} & %%WM_TESTS%% & add/remove, layout, reorder, dimensions, state \\
Persistence & \texttt{persistence.test.js} & %%PERSIST_TESTS%% & save/load, write guard, cleanup \\
InstanceRegistry & \texttt{instance-registry.test.js} & %%IR_TESTS%% & register, prune, CQS, heartbeat \\
Renderer UI & \texttt{renderer-ui.test.js} & %%UI_TESTS%% & DOM, light/DND mode, state sync, rendering \\
\bottomrule
\end{tabularx}
\caption{Tests unitarios por m\'odulo}
\end{table}

\subsection{Tests de Integraci\'on}

\begin{itemize}
    \item \textbf{Concurrencia (z2b.10):} 5 procesos hijo escriben simult\'aneamente al registro, verificando que el file locking previene corrupci\'on. Se ejecuta 3 veces consecutivas para mayor confianza.
    \item \textbf{Pipeline del proceso principal (z2b.15):} Verifica 6 flujos: add-window, focus-change, cleanup, remove-window, layout, y state persistence.
\end{itemize}

\subsection{Resumen de Testing}

\begin{infobox}[Resultado final de tests]
\textbf{%%TOTAL_TESTS%% tests} en \textbf{7 archivos} --- todos pasando en \textbf{$\sim$7 segundos}
\end{infobox}

% === SECTION 6: FASE 3 ===
\section{Fase 3 --- Robustez y Rendimiento}

\subsection{Cache de Callback koffi (z2b.16)}
\texttt{getAvailableWindows()} creaba y destru\'ia un callback FFI en cada invocaci\'on (cada 5 segundos). Se movi\'o la creaci\'on al constructor de \texttt{WindowManager} y la destrucci\'on a un nuevo m\'etodo \texttt{cleanup()}, reduciendo la frecuencia de \texttt{koffi.register/unregister} de 12 veces/minuto a 1 vez por sesi\'on.

\subsection{L\'imite Superior de Dimensiones (z2b.17)}
\texttt{setCustomDimensions()} ahora clampea valores al rango $[200, \min(10000, \text{workArea})]$. El handler IPC y \texttt{onManagedWindowResized} tambi\'en aplican clamping al \'area de trabajo del display actual.

\subsection{DPI Awareness (z2b.18)}
Se declar\'o \texttt{SetProcessDpiAwarenessContext} en \texttt{win32.js} y se invoca con \texttt{DPI\_AWARENESS\_CONTEXT\_PER\_MONITOR\_AWARE\_V2} antes de \texttt{app.whenReady()}. Si la API no est\'a disponible (Windows antiguo), se registra un warning sin interrumpir la ejecuci\'on.

\subsection{Deduplicaci\'on de State Sync (z2b.19)}
Se extrajo la l\'ogica duplicada ($\sim$160 l\'ineas) entre \texttt{refreshManaged()} y el callback \texttt{onStateUpdate} en una \'unica funci\'on \texttt{syncStateToUI(data)} que maneja todos los campos de estado respetando los guards de foco.

% === SECTION 7: REGRESION ===
\section{Regresi\'on Detectada y Corregida}

\begin{regressionbox}[Bug: Lista de ventanas disponibles vac\'ia tras z2b.16]
Al cachear el callback de \texttt{EnumWindows}, el worker omiti\'o la llamada a \texttt{koffi.address()} necesaria para convertir el callback registrado a su direcci\'on de memoria nativa.

\textbf{C\'odigo incorrecto:} \texttt{api.EnumWindows(this.\_enumCallback, 0)}

\textbf{C\'odigo correcto:} \texttt{api.EnumWindows(koffi.address(this.\_enumCallback), 0)}

\textbf{Causa ra\'iz:} El c\'odigo original usaba \texttt{koffi.address(callback)} pero el worker copi\'o solo el nombre de la variable sin el wrapper.

\textbf{Detecci\'on:} Testing manual del usuario.

\textbf{Lecci\'on:} Los tests unitarios con mocks de koffi no detectaron este bug porque el mock no valida la sem\'antica de \texttt{address()} vs paso directo. Se requieren tests E2E o de integraci\'on con el runtime real para validar interacciones FFI.
\end{regressionbox}

% === SECTION 8: METRICAS ===
\section{M\'etricas Cuantitativas}

\begin{table}[H]
\centering
\renewcommand{\arraystretch}{1.5}
\begin{tabularx}{\textwidth}{X r r}
\toprule
\textbf{M\'etrica} & \textbf{Antes} & \textbf{Despu\'es} \\
\midrule
Archivos fuente & 9 & %%METRIC_SRC_FILES%% \\
L\'ineas de c\'odigo (src/) & 3,168 & %%METRIC_SRC_LINES%% \\
Archivos de test & 0 & 7 \\
Tests automatizados & 0 & %%TOTAL_TESTS%% \\
Errores de lint & Sin linter & %%LINT_ERRORS%% \\
Warnings de lint & Sin linter & %%LINT_WARNINGS%% \\
CSP script-src & \texttt{'self' 'unsafe-inline'} & \texttt{'self'} \\
WebPreferences seguridad & 4 propiedades & 6 propiedades \\
File locking (registry) & Ninguno & \texttt{proper-lockfile} \\
Heartbeat de instancias & No exist\'ia & Cada 60 segundos \\
Renderer & 1 archivo (1,715 l\'in.) & 3 archivos \\
Callbacks FFI por minuto & 12 (register+unregister) & 0 (cacheado) \\
DPI awareness & Ninguno & Per-Monitor V2 \\
L\'ogica duplicada (renderer) & $\sim$160 l\'ineas & 0 (syncStateToUI) \\
\bottomrule
\end{tabularx}
\caption{Comparaci\'on cuantitativa antes y despu\'es de la implementaci\'on}
\end{table}

% === SECTION 9: PENDIENTES ===
\section{Tareas Pendientes (Fase 4)}

\begin{table}[H]
\centering
\renewcommand{\arraystretch}{1.3}
\begin{tabularx}{\textwidth}{l c X c}
\toprule
\textbf{Tarea} & \textbf{Prior.} & \textbf{Descripci\'on} & \textbf{Esfuerzo} \\
\midrule
z2b.4 & P2 & Modularizar \texttt{app.js} en componentes UI e IPC & 1 d\'ia \\
z2b.20 & P2 & Configurar \texttt{electron-builder} para installer Windows & 3 d\'ias \\
z2b.21 & P3 & Tests E2E con Playwright para Electron & 3 d\'ias \\
z2b.22 & P3 & Investigar detecci\'on event-driven de ventanas & 2 d\'ias \\
\bottomrule
\end{tabularx}
\caption{Tareas pendientes para la Fase 4}
\end{table}

% === SECTION 10: LECCIONES ===
\section{Lecciones Aprendidas}

\begin{enumerate}
    \item \textbf{Workers paralelos + worktrees = velocidad con costo de merge.} La ejecuci\'on de 4--6 workers simult\'aneos aceler\'o significativamente el desarrollo, pero gener\'o conflictos de merge cuando m\'ultiples workers modificaban archivos compartidos (especialmente \texttt{package.json} y archivos formateados por Prettier).

    \item \textbf{Los mocks de CJS en Vitest requieren t\'ecnicas especiales.} \texttt{vi.mock()} intercepta imports ESM pero no \texttt{require()} de CJS. Los workers resolvieron esto inyectando mocks en el cache de m\'odulos de Node.js o usando \texttt{vi.spyOn()} sobre objetos compartidos.

    \item \textbf{Los tests unitarios con mocks no detectan bugs de integraci\'on FFI.} La regresi\'on de \texttt{koffi.address()} demostr\'o que mockear completamente la capa FFI oculta errores sem\'anticos. Se necesitan tests E2E con el runtime real para validar interacciones nativas.

    \item \textbf{El ordenamiento ``safety-first'' del plan refinado fue acertado.} Implementar file locking (z2b.1) antes de cualquier otro cambio al registry evit\'o introducir bugs de concurrencia durante el desarrollo paralelo.

    \item \textbf{La separaci\'on del renderer es prerequisito de todo lo dem\'as.} Sin la extracci\'on de z2b.3, no se pod\'ia eliminar \texttt{unsafe-inline}, no se pod\'ia hacer linting del JS, y no se pod\'ian escribir tests del renderer.

    \item \textbf{El heartbeat faltante era un bug silencioso cr\'itico.} \texttt{\_startHeartbeat()} se llamaba pero no exist\'ia --- JavaScript no lanza error al llamar \texttt{undefined()} en un try/catch. Sin la auditor\'ia, este bug habr\'ia persistido indefinidamente.
\end{enumerate}

% === SECTION 11: CONCLUSIONES ===
\section{Conclusiones}

La ejecuci\'on del plan refinado transform\'o \textbf{Stack Windows Electron} de un proyecto funcional pero fr\'agil (B+) a una base de c\'odigo robusta, testeable y mantenible (A--). Los cambios no alteraron la experiencia del usuario --- la aplicaci\'on se ve y funciona id\'enticamente --- pero mejoraron fundamentalmente su arquitectura interna.

Los logros m\'as significativos fueron: (1) la eliminaci\'on del riesgo de corrupci\'on de datos en escenarios multi-instancia mediante file locking y CQS, (2) la creaci\'on de una red de seguridad de %%TOTAL_TESTS%% tests automatizados que cubren la l\'ogica cr\'itica del negocio, y (3) el endurecimiento de seguridad que cerr\'o el vector de XSS m\'as com\'un en aplicaciones Electron.

El proyecto queda en una posici\'on s\'olida para abordar la Fase 4 (distribuci\'on y testing E2E), con la confianza de que los cambios futuros ser\'an validados autom\'aticamente por la infraestructura de testing establecida.

\end{document}
